%% kapa: exam=1 round=28 year=2017 period=2 form=A subject=accounting subject_ko=회계학 seq=5 number=45
%% source: Q-Net 감정평가사 기출문제 (cst00309) artlSeq=5199603 file=extracted/1008/5199603/2017년 제28회 감정평가사 1차시험 문제지/2교시 A형.pdf
%% conversion: kordoc 4.24.51 -> tools/gen1.py (auto); review=proofread-sonnet; flags=none
\begin{question}{45}
보고기간후사건에 관한 설명으로 옳지 않은 것은?
\begin{choices}
  \item 보고기간 후부터 재무제표 발행승인일 전 사이에 배당을 선언한 경우에는 보고기간말에 부채로 인식한다.
  \item 보고기간말 이전에 구입한 자산의 취득원가나 매각한 자산의 대가를 보고기간 후에 결정하는 경우는 수정을 요하는 보고기간후사건이다.
  \item 보고기간말과 재무제표 발행승인일 사이에 투자자산의 공정가치의 하락은 수정을 요하지 않는 보고기간후사건이다.
  \item 보고기간 후에 발생한 화재로 인한 주요 생산 설비의 파손은 수정을 요하지 않는 보고기간후사건이다.
  \item 경영진이 보고기간 후에, 기업을 청산하거나 경영활동을 중단할 의도를 가지고 있다고 판단하는 경우에는 계속기업의 기준에 따라 재무제표를 작성해서는 아니 된다.
\end{choices}
\end{question}
\answer{1}
